%-------------------------------------------------------------------------------
% SECTION TITLE
%-------------------------------------------------------------------------------
\cvsection{Projects}


%-------------------------------------------------------------------------------
% CONTENT
%-------------------------------------------------------------------------------
\begin{cventries}

%---------------------------------------------------------
  \cventry
    {Founder}
    {GrowYumi}
    {Personal WeChat Mini Program}
    {Feb. 2022 - Present}
    {
      \begin{cvitems}
        \item {Built and operated a personal WeChat Mini Program for couples, reaching \textbf{\textcolor{awesome-red}{100,000+}} registered users.}
        \item {With access to the WeChat advertising platform, monthly income can reach RMB 1,000 during peak periods.}
      \end{cvitems}
    }

%---------------------------------------------------------
  \cventry
    {Personal Project}
    {A Simplified vLLM Core}
    {LLM Inference Systems}
    {Dec. 2025 - Present}
    {
      \begin{cvitems}
        \item {Rebuilt core vLLM-style inference components from scratch in PyTorch/CUDA, including paged KV-cache block management, continuous batching, and a lightweight LLM serving engine.}
        \item {Implemented a FlashAttention-style prefill CUDA kernel and benchmarked it against dense baselines, achieving about 10\% lower prefill latency than a naive CUDA attention kernel (4.64 ms vs. 5.15 ms).}
        \item {Designed KV-cache writeback and INT8 KV-cache quantization, reducing KV memory usage by about 50\% (0.0078 GB to 0.0039 GB) with low quantization overhead (0.052 ms quantize, 0.145 ms dequantize).}
      \end{cvitems}
    }

%---------------------------------------------------------
  \cventry
    {Graduation Project}
    {RGB-Infrared Fusion for Vehicle-Scene Segmentation}
    {Sichuan University}
    {Sep. 2020 - Nov. 2021}
    {
      \begin{cvitems}
        \item {Improved semantic segmentation under low-light and glare conditions with RGB-infrared fusion.}
        \item {Used PyTorch and a temporal-difference method to improve recognition efficiency.}
      \end{cvitems}
    }

%---------------------------------------------------------
\end{cventries}
